\clearpage
\phantomsection% 
\addcontentsline{toc}{chapter}{ABSTRACT}
\thispagestyle{fancy}

\begin{center}
	\large \bfseries \MakeUppercase{Abstract}\\
%	\normalsize \normalfont {\thetitleEN}\\
%	\normalsize \normalfont {\theauthor}\\
	\bigskip
	
	\normalsize \normalfont \justifying \singlespacing
	The advancement of generative artificial intelligence has increased the ability to create highly realistic \textit{deepfake} videos, raising challenges for information security and digital trust. \textit{Deepfake} detection in real-world video data is increasingly difficult due to variations in visual quality, pose, illumination, resolution, and temporal characteristics. This study aims to compare a spatial approach based on Vision Transformer (ViT) and a spatiotemporal approach based on VideoMAE for \textit{deepfake} video classification. \textit{WildDeepfake} is used as the main dataset, while \textit{Celeb-DF v2} is used as an external dataset for \textit{zero-shot cross-dataset} evaluation. The research stages include data preprocessing, 46-\textit{frame} clip formation, \textit{strided temporal sampling}, model training, training configuration evaluation, spatial and temporal input configuration evaluation, and model assessment using accuracy, precision, recall, F1-score, AUC-ROC, and confusion matrix. In the baseline experiment on \textit{WildDeepfake}, the spatiotemporal model achieved 85.8\% accuracy, 84.5\% F1-score, and 92.3\% AUC, while the spatial model achieved 79.2\% accuracy, 76.7\% F1-score, and 86.8\% AUC. These results indicate that modeling inter-\textit{frame} relationships using VideoMAE contributes to improved performance in internal evaluation. However, in the cross-dataset evaluation on \textit{Celeb-DF v2}, the spatial model achieved better performance with 72.9\% accuracy and 77.7\% AUC, while the spatiotemporal model achieved 64.2\% accuracy and 72.6\% AUC. Therefore, the spatiotemporal approach is more effective in the internal \textit{WildDeepfake} evaluation, whereas the spatial approach demonstrates better generalization in the cross-dataset scenario.\par
    
    \vspace{20pt}
	\raggedright\textbf{Keywords}: \textit{deepfake}, ViT, VideoMAE, spatiotemporal, \textit{WildDeepfake}
	
	\vfill
	
\end{center}
\clearpage
